\begin{tikzpicture}[
  x=1cm,y=1cm,
  font=\small,
  class/.style={draw,rounded corners=2pt,text width=34mm,minimum height=13mm,align=center,fill=blue!7,inner sep=3pt},
  label/.style={class,fill=orange!11},
  process/.style={class,fill=green!10},
  artifact/.style={class,fill=purple!8},
  relation/.style={-{Stealth[length=2mm]},thick},
  isa/.style={relation,dashed},
  edgeword/.style={font=\footnotesize,fill=white,inner sep=1.5pt},
  lane/.style={font=\small\bfseries,text=black!70,anchor=west}
]
% La mayor separación vertical y la tipografía ampliada priorizan la lectura
% al ancho de una página A4 de una columna.
\node[lane] at (-1.75,1.05) {Evidencia y etiquetado};
\node[class] (video) at (0,0) {Video público};
\node[class] (subtitle) at (6,0) {Pista de subtítulos};
\node[class] (chunk) at (12,0) {Chunk de evidencia};
\node[label] (labeling) at (18,0) {Etiquetado};
\draw[relation] (video.east)--node[edgeword,above]{contiene}(subtitle.west);
\draw[relation] (subtitle.east)--node[edgeword,above]{segmenta}(chunk.west);
\draw[relation] (chunk.east)--node[edgeword,above]{etiqueta}(labeling.west);

\node[artifact] (dataset) at (12,-2.9) {Artefacto de datos\\versión y partición};
\node[label] (concept) at (18,-2.9) {Concepto de\\moderación};
\draw[relation] (chunk.south)--node[edgeword,right]{integra}(dataset.north);
\draw[relation] (labeling.south)--node[edgeword,right]{clasifica}(concept.north);

\node[lane] at (-1.75,-4.7) {Taxonomía operativa};
\node[label] (safe) at (6,-6.0) {Estado \texttt{SEGURO}\\derivado};
\node[label] (damage) at (12,-6.0) {Categoría operativa\\multietiqueta};
\node[label] (fine) at (18,-6.0) {Fenómeno fino\\o señal de revisión};
\coordinate (conceptdrop) at (18,-4.4);
\coordinate (conceptbus) at (15,-4.4);
\draw[dashed,thick] (concept.south) -- (conceptdrop) -- (conceptbus);
\fill (conceptbus) circle (1.1pt);
\draw[isa] (conceptbus) -| (damage.north);
\draw[isa] (conceptbus) -| (fine.north);
\draw[relation] (damage.west)--node[edgeword,above,align=center]{si ninguna\\se activa}(safe.east);

\node[lane] at (-1.75,-8.2) {Inferencia, decisión y retroalimentación};
\node[process] (model) at (0,-9.6) {Modelo versionado};
\node[process] (prediction) at (6,-9.6) {Predicción\\score y umbral};
\node[process] (review) at (12,-9.6) {Decisión humana\\aceptar, modificar\\o rechazar};
\node[artifact] (feedback) at (18,-9.6) {Registro reentrenable\\decisiones versionadas};
\draw[relation] (model.east)--node[edgeword,above]{produce}(prediction.west);
\draw[relation] (prediction.east)--node[edgeword,above,align=center]{si requiere\\revisión}(review.west);
\draw[relation] (review.east)--node[edgeword,above]{registra}(feedback.west);

% Carril izquierdo: solo los datos de entrenamiento alimentan el modelo.
\coordinate (trainlane) at (-2.35,-2.9);
\draw[relation] (dataset.west) -- (trainlane) |- (model.west);
\node[edgeword,anchor=east] at (-2.35,-7.35) {entrena};

% Carril central: la predicción refiere a una categoría sin atravesar nodos.
\coordinate (predictlane) at (9,-7.65);
\draw[relation] (prediction.north) -- (prediction.north |- predictlane) -| (damage.south);
\node[edgeword] at (9,-7.65) {predice};

% La retroalimentación crea una versión futura; no modifica el corpus evaluado.
\coordinate (futurelane) at (18,-11.35);
\draw[isa] (feedback.south) -- (futurelane) -| (model.south);
\node[edgeword,below] at (9,-11.35) {futura iteración DSR};
\end{tikzpicture}
